\section{Reweighting Pipeline}

% ------------------------------------------------------------------
\subsection{ResiDual: Implementation Details}
\label{app:residual_impl}
% ------------------------------------------------------------------
 
This section describes the full implementation of ResiDual (RD), 
a learnable adaptation of the residual reweighting idea integrated 
into a CLAP-based model. The pipeline consists of two sequential phases: a 
\emph{fitting pass} that estimates a fixed PCA basis for each attention head 
(\S\ref{app:fitting}), followed by an \emph{optimisation phase} in which 
the per-head weight vectors are 
further refined as free parameters via gradient descent to maximise zero-shot 
classification accuracy (\S\ref{app:optim}). The forward-pass reweighting 
that uses these learned weights at inference time is described in \S\ref{app:forward_rw}.

 
\paragraph{Module injection via forward hooks.}
RD injects its reweighting through PyTorch forward
hooks. For each target block $(\ell, b)$ with $\ell \in \mathcal{L}$, a
dedicated \texttt{AttentionHook} object is registered on the corresponding
\texttt{WindowAttention} module immediately before the forward pass begins
and removed as soon as the forward pass for that layer completes.

The hook intercepts the tuple $(\mathbf{x}_{\mathrm{out}},\,\mathbf{A})$
returned by \texttt{WindowAttention}, where $\mathbf{x}_{\mathrm{out}} \in
\mathbb{R}^{B' \times N \times C}$ is the projected attention output and
$\mathbf{A}$ contains the attention weights. It then reconstructs the
per-head pre-projection tensors $\mathbf{H}_{\ell,b,h}$ by re-applying
the frozen QKV projection to the hook input, isolating the value slice $\mathbf{V}_h$,
and computing $\mathbf{x}_{\mathrm{heads}} = \mathbf{A}\,\mathbf{V}$. This
yields tensors of shape $(B', H_\ell, N, d_h)$ without any additional
trainable operation. No structural modification of the original model is
performed: all pretrained parameters — the fused QKV projection
$W^{QKV}_{\ell,b}$, the output projection $W^O_{\ell,b}$ with bias
$\mathbf{b}^O_{\ell,b}$, and the relative position bias table — remain
exactly as loaded from the checkpoint. The only additions are the learnable
weight vectors $\boldsymbol{\lambda}_{\ell,b,h}$ and their associated
PCA buffers, stored in a separate \texttt{nn.ModuleDict} that is indexed by
$(\ell, b, h)$.
 
% ------------------------------------------------------------------
\subsubsection{Fitting Pass: PCA Basis Estimation}
\label{app:fitting}
% ------------------------------------------------------------------
 
\paragraph{Data collection.}
The fitting pass is run with gradients disabled and the model in evaluation
mode. For each audio clip (up to $N_{\max}$ clips), a forward pass is
executed with data collection enabled. The hook registered at each target
block intercepts, for every head $h$, the pre-reweighting tensor
$\mathbf{H}_{\ell,b,h} \in \mathbb{R}^{B' \times H_\ell \times N \times d_h}$
and flattens it along the spatial and batch axes to obtain
$\mathbb{R}^{(B' \cdot N) \times d_h}$ before appending it to an
accumulation buffer for $(\ell, b, h)$. Here $B'$ is the number of
attention windows times the batch size and $N = M = 64$ is the number of
tokens per window. After all clips have been processed, the buffers are
concatenated to form the fitting matrix
\begin{equation}
    \widetilde{\mathbf{R}}_{\ell,b,h}
    \;\in\; \mathbb{R}^{N_{\mathrm{fit}} \times d_h},
    \qquad
    N_{\mathrm{fit}} \;=\; \sum_{\mathrm{batches}} B \cdot N_w^\ell \cdot M,
    \label{eq:fitting_matrix}
\end{equation}
where the summation is over all processed batches of size $B$.
 
\paragraph{PCA decomposition.}
Standard PCA is applied to $\widetilde{\mathbf{R}}_{\ell,b,h}$: the
empirical mean $\boldsymbol{\mu}_{\ell,b,h} \in \mathbb{R}^{d_h}$ is
computed and subtracted, and the empirical covariance is decomposed as
\begin{align}
    \mathbf{C}_{\ell,b,h}
    &= \frac{1}{N_{\mathrm{fit}} - 1}
       \Bigl(\widetilde{\mathbf{R}}_{\ell,b,h}
             - \mathbf{1}\boldsymbol{\mu}_{\ell,b,h}^\top\Bigr)^\top
       \Bigl(\widetilde{\mathbf{R}}_{\ell,b,h}
             - \mathbf{1}\boldsymbol{\mu}_{\ell,b,h}^\top\Bigr) \notag \\
    &= \boldsymbol{\Phi}_{\ell,b,h}\,
       \operatorname{diag}(\lambda^{\mathrm{ev}}_1, \ldots, \lambda^{\mathrm{ev}}_{d_h})\,
       \boldsymbol{\Phi}_{\ell,b,h}^\top,
    \quad
    \lambda^{\mathrm{ev}}_1 \ge \cdots \ge \lambda^{\mathrm{ev}}_{d_h} \ge 0,
    \label{eq:pca_decomp}
\end{align}
where $\boldsymbol{\Phi}_{\ell,b,h} \in \mathbb{R}^{d_h \times d_h}$ is the
orthogonal matrix of eigenvectors and $\lambda^{\mathrm{ev}}_j$ denotes the
$j$-th eigenvalue of the empirical covariance (we use the superscript
``ev'' to distinguish these fixed eigenvalues from the trainable weight
vector $\boldsymbol{\lambda}_{\ell,b,h}$ introduced below).
The full basis $\boldsymbol{\Phi}_{\ell,b,h}$ and the mean
$\boldsymbol{\mu}_{\ell,b,h}$ are stored as non-trainable buffers; the
temporary accumulation buffers are then cleared to free memory.
 
\paragraph{Task-relevant subspace and initialisation of $\boldsymbol{\lambda}$.}
For each head, the index
\begin{equation}
    k_{\ell,b,h}
    \;=\; \min\!\left\{
        j \in \{1,\ldots,d_h\} :
        \sum_{i=1}^{j} \frac{\lambda^{\mathrm{ev}}_i}{\sum_{i'} \lambda^{\mathrm{ev}}_{i'}}
        \;\ge\; \rho_{\mathrm{var}}
    \right\}
    \label{eq:k_threshold}
\end{equation}
identifies the smallest number of principal components that explain at least
a fraction $\rho_{\mathrm{var}}$ of the total variance, providing a
per-head estimate of the effective dimensionality of the task-relevant
subspace. The learnable weight vector $\boldsymbol{\lambda}_{\ell,b,h} \in
\mathbb{R}^{d_h}$ is then initialised as
\begin{equation}
    \lambda_j^{(0)}
    \;=\; \begin{cases}
        1 & j \le k_{\ell,b,h}, \\
        0 & j > k_{\ell,b,h}.
    \end{cases}
    \label{eq:lambda_init}
\end{equation}
This initialisation corresponds to the identity map restricted to the
task-relevant subspace and a complete suppression of the noise components:
before any optimisation, RD$^*$\xspace acts as a projection-and-reconstruction
operator that retains the top-$k$ PCA directions unchanged and zeros out
the remainder.
 
% ------------------------------------------------------------------
\subsubsection{Forward Pass with Spectral Reweighting}
\label{app:forward_rw}
% ------------------------------------------------------------------
 
At inference time, the hook at block $(\ell, b)$ intercepts the per-head
tensors $\mathbf{H}_{\ell,b,h} \in \mathbb{R}^{B' \times N \times d_h}$
and applies the following sequence of operations for each head $h$.
 
\textbf{Step 1 — Centering.}
\begin{equation}
    \widetilde{\mathbf{H}}_{\ell,b,h}
    \;=\; \mathbf{H}_{\ell,b,h} - \boldsymbol{\mu}_{\ell,b,h}
    \;\in\; \mathbb{R}^{B' \times N \times d_h},
\end{equation}
where $\boldsymbol{\mu}_{\ell,b,h}$ is broadcast over the first two axes.
 
\textbf{Step 2 — Projection onto the full PCA basis.}
\begin{equation}
    \mathbf{Z}_{\ell,b,h}
    \;=\; \widetilde{\mathbf{H}}_{\ell,b,h}\,\boldsymbol{\Phi}_{\ell,b,h}
    \;\in\; \mathbb{R}^{B' \times N \times d_h}.
\end{equation}
Unlike fixed-weight variants that truncate to $K < d_h$ components, RD$^*$\xspace
projects onto the \emph{full} basis $\boldsymbol{\Phi}_{\ell,b,h} \in
\mathbb{R}^{d_h \times d_h}$, delegating the effective truncation to the
learned weights.
 
\textbf{Step 3 — Per-component rescaling by $\boldsymbol{\lambda}_{\ell,b,h}$.}
\begin{equation}
    \widehat{\mathbf{Z}}_{\ell,b,h}
    \;=\; \mathbf{Z}_{\ell,b,h} \odot \boldsymbol{\lambda}_{\ell,b,h}^\top
    \;\in\; \mathbb{R}^{B' \times N \times d_h},
    \label{eq:rd_reweight}
\end{equation}
where $\boldsymbol{\lambda}_{\ell,b,h}^\top$ is broadcast over the first two
axes. Because $\boldsymbol{\lambda}_{\ell,b,h}$ is a free real-valued
parameter, each entry $\lambda_j$ can amplify ($|\lambda_j| > 1$), attenuate
($|\lambda_j| \in (0,1)$), suppress ($\lambda_j = 0$), or invert
($\lambda_j < 0$) the corresponding principal component.
 
\textbf{Step 4 — Reconstruction in the original space.}
\begin{equation}
    \mathbf{H}^*_{\ell,b,h}
    \;=\; \widehat{\mathbf{Z}}_{\ell,b,h}\,\boldsymbol{\Phi}_{\ell,b,h}^\top
    \;\in\; \mathbb{R}^{B' \times N \times d_h}.
    \label{eq:reweighted_head}
\end{equation}
The empirical mean is \emph{not} restored after reconstruction. Since
$\mathbf{H}^*_{\ell,b,h}$ is immediately passed to the output projection
$W^O_{\ell,b}$, which has its own learnable bias $\mathbf{b}^O_{\ell,b}$,
re-adding $\boldsymbol{\mu}_{\ell,b,h}$ would introduce a systematic,
data-dependent shift in the residual stream that is already absorbed by
$\mathbf{b}^O_{\ell,b}$.
 
Combining Steps~1--4, the full RD$^*$\xspace transform can be written compactly as
\begin{equation}
    \mathrm{RD}^*(\mathbf{H}_{\ell,b,h},\,\boldsymbol{\lambda}_{\ell,b,h})
    \;=\;
    \boldsymbol{\Phi}_{\ell,b,h}\,
    \operatorname{diag}(\boldsymbol{\lambda}_{\ell,b,h})\,
    \boldsymbol{\Phi}_{\ell,b,h}^\top
    \bigl(\mathbf{H}_{\ell,b,h} - \boldsymbol{\mu}_{\ell,b,h}\bigr).
    \label{eq:rd_star}
\end{equation}
The tensor $\mathbf{H}^*_{\ell,b,h}$ replaces $\mathbf{H}_{\ell,b,h}$ in
Eq.~\eqref{eq:wmsa_out}, so the output projection $W^O_{\ell,b}$ and the
residual-stream update of Eq.~\eqref{eq:residual_attn} proceed unchanged.
The gradient of any downstream loss flows through
$\mathrm{RD}^*$ and updates only $\boldsymbol{\lambda}_{\ell,b,h}$; the
basis $\boldsymbol{\Phi}_{\ell,b,h}$, the mean $\boldsymbol{\mu}_{\ell,b,h}$,
and all original model parameters are held fixed.
 
% ------------------------------------------------------------------
\subsubsection{Optimisation of $\boldsymbol{\lambda}$ via Zero-Shot Accuracy}
\label{app:optim}
% ------------------------------------------------------------------
 
After the fitting pass has fixed the PCA bases and initialised
$\boldsymbol{\lambda}$ via Eq.~\eqref{eq:lambda_init}, the second phase
optimises the weight vectors to maximise zero-shot classification accuracy
on a held-out validation set.
 
\paragraph{Parameter isolation.}
All model parameters except the weight vectors
$\{\boldsymbol{\lambda}_{\ell,b,h}\}_{\ell \in \mathcal{L},\,b,\,h}$
are frozen by setting \texttt{requires\_grad = False}. This ensures that
the only degrees of freedom during optimisation are the $|\mathcal{L}|
\times \sum_{\ell \in \mathcal{L}} B_\ell \cdot H_\ell$ scalar values
that control the per-component gain of each attention head in the target
stages. The original CLAP audio projection head $P : \mathbb{R}^{D_3}
\to \mathbb{R}^{d_{\mathrm{proj}}}$ is kept frozen but is included in
the computational graph so that gradients flow from the similarity space
back to $\boldsymbol{\lambda}$.
 
\paragraph{Objective.}
Given a batch of audio clips with ground-truth class labels
$y \in \{1,\ldots,C\}$ and pre-computed, $\ell_2$-normalised text
embeddings $\mathbf{t}_c \in \mathbb{R}^{d_{\mathrm{proj}}}$ for each
class $c$, the optimisation objective is the cross-entropy loss over
zero-shot logits:
\begin{equation}
    \mathcal{J}(\boldsymbol{\lambda})
    \;=\;
    -\,\mathbb{E}_{(x,y)}\!\left[
        \log \frac{\exp\bigl(\hat{\mathbf{a}}(x)^\top \mathbf{t}_y\bigr)}
                  {\sum_{c=1}^{C} \exp\bigl(\hat{\mathbf{a}}(x)^\top \mathbf{t}_c\bigr)}
    \right],
    \label{eq:optim_obj}
\end{equation}
where $\hat{\mathbf{a}}(x) = \ell_2\text{-norm}\bigl(P(\mathrm{latent}(x))\bigr) \in
\mathbb{R}^{d_{\mathrm{proj}}}$ is the $\ell_2$-normalised audio
embedding of clip $x$ and $\mathrm{latent}(x)$ is the spatially averaged
final-stage output of HTS-AT after applying the RD$^*$\xspace transforms.
Minimising $\mathcal{J}$ is equivalent to maximising zero-shot accuracy
in expectation.
 
\paragraph{Optimiser and schedule.}
The weight vectors are updated with Schedule-Free Adam~\cite{defazio2024schedulefree}
at learning rate $\eta$, which eliminates the need for a separate learning-rate
schedule while matching or exceeding the performance of Adam with cosine
decay on this type of fine-tuning task. When the \texttt{schedulefree}
package is unavailable, standard Adam is used as a fallback.
 
\paragraph{Early stopping and best-$\boldsymbol{\lambda}$ restoration.}
At the end of each epoch, zero-shot accuracy on the validation set is
computed. A snapshot of the current $\boldsymbol{\lambda}$ vectors is saved
whenever a new accuracy maximum is achieved. Training halts early if no
improvement is observed for $p$ consecutive epochs. At termination, the
best snapshot is restored, guaranteeing that the final weight vectors
correspond to the configuration that achieved the highest observed
validation accuracy rather than to the last iterate.
 
\paragraph{Hyperparameter summary.}
Table~\ref{tab:residual_hparams} summarises all hyperparameters introduced
by RD$^*$\xspace and the default values used in all experiments reported in
this paper.
 
\begin{table}[h]
\centering
\caption{RD$^*$ hyperparameters.}
\label{tab:residual_hparams}
\begin{tabular}{lll}
\toprule
\textbf{Symbol} & \textbf{Description} & \textbf{Default} \\
\midrule
$\mathcal{L}$
    & Target stage indices for reweighting
    & $\{1, 2, 3\}$ \\
$\rho_{\mathrm{var}}$
    & Variance threshold for $k$ estimation (Eq.~\eqref{eq:k_threshold})
    & $0.95$ \\
$\eta$
    & Learning rate for Schedule-Free Adam
    & $10^{-2}$ \\
$p$
    & Early-stopping patience (epochs)
    & $5$ \\
$T_{\max}$
    & Maximum number of optimisation epochs
    & $30$ \\
$N_{\max}$
    & Maximum fitting samples
    & $10\,000$ \\
\bottomrule
\end{tabular}
\end{table}
 
% ------------------------------------------------------------------
\paragraph{Hook-based extraction (updated note).}
% ------------------------------------------------------------------
 
The hook-based mechanism described in \S\ref{sec:aggregation} for the
dimensionality analysis is also the mechanism through which RD$^*$\xspace applies
reweighting during inference. The same \texttt{AttentionHook} object
operates in two distinct modes depending on the \texttt{collect\_for\_fitting}
flag: in collection mode it accumulates raw head activations for PCA
estimation without modifying the forward pass output; in reweighting mode
it replaces the projected attention output with the RD$^*$\xspace-transformed
version before returning it to the residual stream. Crucially, in both
modes the hooks are \emph{temporary}: they are installed at the start of
each call to \texttt{\_forward\_layer\_with\_attention\_hooks} and removed
immediately upon its completion, so the underlying \texttt{WindowAttention}
modules are never structurally modified.